Qwen's paired alignment gain is 0.039, with 95\% CI [0.016, 0.062]. The gain interval lies above zero. Its below-mean selection rate is 36.3\%, compared with matched random probability 43.1\%. A positive BMR alone is not evidence of degradation relative to random. The unanimity rule changes the evaluated subset from 300 majority-policy prompts to 121 prompts. Regret point estimates are 0.135 and 0.074, respectively, but these means concern different supports. On unanimous prompts, the original-order choice is unchanged. The retained-subset gain checks selection value after filtering; it is not an effect of changing those choices. Selection also raises Qwen's mean stereotype rating relative to matched random selection by 0.023 (pointwise 95\% CI [0.002, 0.045]), while lowering missing-expectation ratings. The stereotype difference is exploratory, but it shows that an alignment gain need not carry over to every cultural outcome. Cross-model agreement has paired alignment gain -0.060 (95\% CI [-0.111, -0.011]) on its accepted prompts. In this configuration, agreement with a weaker judge that shares a slot preference does not identify better-than-random choices; a stronger or position-debiased second judge could behave differently.
